\documentclass[a4paper]{arrowhead}

\usepackage[yyyymmdd]{datetime}
\usepackage{etoolbox}
\usepackage[utf8]{inputenc}
\usepackage[T1]{fontenc}
\usepackage{multirow}
\usepackage{longtable}
\usepackage{tikz}
\usepackage{adjustbox}
\AtBeginEnvironment{longtable}{\small}
\BeforeBeginEnvironment{tikzpicture}{\begin{adjustbox}{max width=\textwidth}}
\AfterEndEnvironment{tikzpicture}{\end{adjustbox}}
\usetikzlibrary{positioning,shapes.multipart,arrows.meta}

\renewcommand{\dateseparator}{-}
\setlength{\emergencystretch}{3em}
\lstset{language=HTTP,breakatwhitespace=false,columns=fullflexible}

\newcommand{\fparam}[1]{\textit{\textcolor{ArrowheadBlue}{#1}}}

%% Special references
\newcommand{\mref}[2][]{{\textcolor{ArrowheadPurple}{\hyperref[sec:model:#2]{\ifstrempty{#1}{#2}{#1}}}}}
\newcommand{\oref}[1]{{\textcolor{ArrowheadBlue}{\hyperref[sec:operations:#1]{#1}}}}
\newcommand{\pdef}[1]{{\textcolor{ArrowheadGrey}{#1 \label{sec:model:primitives:#1} \label{sec:model:primitives:#1s}}}}
\newcommand{\pref}[1]{{\textcolor{ArrowheadGrey}{\hyperref[sec:model:primitives:#1]{#1}}}}

% as template_idd.tex; the input and output arguments are passed already formatted (\mref or \pref)
\newrobustcmd\osubsection[5]{
  \addtocounter{subsection}{1}
  \addcontentsline{toc}{subsection}{\protect\numberline{\thesubsection}operation \textcolor{ArrowheadBlue}{#1}}
  \subsection*{
    \thesubsection\quad
    #2 \textcolor{ArrowheadPurple}{#3} \\
    \small
    \hspace*{0.075\textwidth}\begin{minipage}{0.1\textwidth}
      \vspace*{1mm}
      Operation: \\
      \notblank{#4}{Input: \\}{}
      \notblank{#5}{Output: \\}{}
    \end{minipage}%
    \begin{minipage}{0.825\textwidth}
      \vspace*{1mm}
      \textcolor{ArrowheadBlue}{#1} \\
      \notblank{#4}{#4 \\}{}
      \notblank{#5}{#5 \\}{}
    \end{minipage}
  }
  \label{sec:operations:#1}
}
\newrobustcmd\msubsection[2]{
  \addtocounter{subsection}{1}
  \addcontentsline{toc}{subsection}{\protect\numberline{\thesubsection}#1 \textcolor{ArrowheadPurple}{#2}}
  \subsection*{\thesubsection\quad#1 \textcolor{ArrowheadPurple}{#2}}
  \label{sec:model:#2} \label{sec:model:#2s}
}
%%

\begin{document}
\sloppy

%% Arrowhead Document Properties
\ArrowheadTitle{plantDescriptionManagement HTTP/TLS/JSON}
\ArrowheadServiceID{plantDescriptionManagement}
\ArrowheadType{Interface Design Description}
\ArrowheadTypeShort{IDD}
\ArrowheadVersion{5.2.2}
\ArrowheadDate{2026-10-30}
\ArrowheadAuthor{Jerker Delsing}
\ArrowheadStatus{RELEASE CANDIDATE}
\ArrowheadContact{jerker.delsing@ltu.se}
\ArrowheadFooter{\href{www.arrowhead.eu}{www.arrowhead.eu}}
\ArrowheadSetup
%%

%% Front Page
\begin{center}
  \vspace*{1cm}
  \huge{\arrowtitle}

  \vspace*{0.2cm}
  \LARGE{\arrowtype}
  \vspace*{1cm}

  \Large{Service ID: \textit{"\arrowid"}}

  \vspace*{\fill}

  % Front Page Image

  \vspace*{1cm}
  \vspace*{\fill}

  % Front Page Abstract
  \begin{abstract}
    This document describes the HTTP/1.1 protocol with TLS payload security and JSON payload encoding variant of the \texttt{plant\allowbreak{}Description\allowbreak{}Management} service of the PlantDescriptionEngine system.
  \end{abstract}

  \vspace*{1cm}
\end{center}

\newpage
%%

%% Table of Contents
\tableofcontents
\newpage
%%

\section{Overview}
\label{sec:overview}

This document describes the \texttt{plant\allowbreak{}Description\allowbreak{}Management} service interface of the PlantDescriptionEngine system, which provides management of plant descriptions: the engineered systems of the plant, the services they consume and produce (ports), and the connections between ports.

It's implemented using protocol, encoding as stated in below table:
\begin{table}[ht!]
  \centering
  \begin{tabular}{|l|l|l|}
    \hline
    \rowcolor{gray!33} Profile type & Type & Version \\ \hline
    Transfer protocol & HTTP & 1.1 \cite{rfc9112} \\ \hline
    Data encryption & TLS & 1.3 \cite{rfc8446}, 1.2 \cite{rfc5246} \\ \hline
    Encoding & JSON & RFC 8259 \cite{rfc8259} \\ \hline
    Compression & N/A & - \\ \hline
    Semantics & N/A & - \\ \hline
    Ontology & N/A & - \\ \hline
  \end{tabular}
  \caption{Communication and semantics details used for the \texttt{plant\allowbreak{}Description\allowbreak{}Management} service interface}
  \label{tab:comunication_semantics_profile}
\end{table}

This document provides the Interface Design Description IDD to the \textit{plantDescriptionManagement -- Service Description} document \cite{pde-plantDescriptionManagement-sd}.
For further details about how this service is meant to be used, please consult that document.

The rest of this document describes how to realize the \texttt{plant\allowbreak{}Description\allowbreak{}Management} service
interface in detail. Both in terms of its operations (Section \ref{sec:operations}) and its information model (Section \ref{sec:model}).

\subsection{Interface registration}

The service is registered in the ServiceRegistry with the interface template \texttt{generic\_\allowbreak{}https}, policy \texttt{CERT\_\allowbreak{}AUTH}, access address \texttt{plantdescriptionengine}, port 8432 and base path \texttt{/\allowbreak{}plantdescriptionengine/\allowbreak{}mgmt}. The \texttt{operations} property of the interface lists the operations below with their method and path relative to the base path (path parameters are placeholders such as \texttt{\{id\}} that the consumer fills in):

\noindent\begin{longtable}{| p{5.2cm} | p{1.6cm} | p{8.4cm} |} \hline
\rowcolor{gray!33} Operation & Method & Path \\ \hline \endhead
\texttt{query} & GET & \texttt{/\allowbreak{}plant-\allowbreak{}descriptions} \\ \hline
\texttt{get} & GET & \texttt{/\allowbreak{}plant-\allowbreak{}descriptions/\allowbreak{}\{id\}} \\ \hline
\texttt{create} & POST & \texttt{/\allowbreak{}plant-\allowbreak{}descriptions} \\ \hline
\texttt{replace} & PUT & \texttt{/\allowbreak{}plant-\allowbreak{}descriptions/\allowbreak{}\{id\}} \\ \hline
\texttt{update} & PATCH & \texttt{/\allowbreak{}plant-\allowbreak{}descriptions/\allowbreak{}\{id\}} \\ \hline
\texttt{remove} & DELETE & \texttt{/\allowbreak{}plant-\allowbreak{}descriptions/\allowbreak{}\{id\}} \\ \hline
\end{longtable}

\newpage

\section{Service Operations}
\label{sec:operations}

The interfaces of the \texttt{plant\allowbreak{}Description\allowbreak{}Management} service, its operations, data models and implementation are provided below.
In particular, each subsection first names the HTTP method and path used to call the operation, after which it names the abstract operation from the \texttt{plant\allowbreak{}Description\allowbreak{}Management} SD document, as well as input and output types.

All operations require the requester's identity, presented according to the provider's authentication policy: the client certificate (mutual TLS) and, with the \texttt{outsourced} policy, \texttt{Authorization: Bearer IDENTITY-TOKEN//<token>}; systems holding the Authentication secret key send \texttt{Bearer AUTHENTICATOR-KEY//<system>//<hash>}, and with the \texttt{declared} policy \texttt{Bearer SYSTEM//<system>} is sent. Request and response bodies are \texttt{application/json} unless stated otherwise.
Errors are answered with an \mref{ErrorMessage} body and the status codes \texttt{400 Bad Request} (invalid parameter), \texttt{401 Unauthorized} (missing or invalid identity), \texttt{403 Forbidden} (not permitted, e.g. by the management, authorization or blacklist filter), \texttt{404 Not Found}, \texttt{408 Request Timeout}, \texttt{423 Locked}, \texttt{500 Internal Server Error} and \texttt{503 Service Unavailable} (a system it depends on failed).

\osubsection{query}{GET}{/\allowbreak{}plantdescriptionengine/\allowbreak{}mgmt/\allowbreak{}plant-descriptions?\allowbreak{}page=\fparam{\{page\}}\allowbreak{}\&size=\fparam{\{size\}}\allowbreak{}\&direction=\fparam{\{direction\}}\allowbreak{}\&sort\_field=\{sort\_field\}\allowbreak{}\&active=\fparam{\{active\}}}{}{\mref[Plant\allowbreak{}Description\allowbreak{}List\allowbreak{}Response]{PlantDescriptionListResponse}}

Returns a page of plant descriptions, optionally only the active one.

Query parameters: \fparam{page} (Integer, optional); \fparam{size} (Integer, optional); \fparam{direction} (String, optional); \fparam{sort\_field} (String, optional); \fparam{active} (Boolean, optional). List values are repeated or comma-separated.

A successful call returns \texttt{200 OK} with the output in the body.

\begin{lstlisting}[label={lst:query},caption={An \oref{query} invocation.}]
GET /plantdescriptionengine/mgmt/plant-descriptions?page=<page>&size=<size>&direction=<direction>&sort_field=<sort_field>&active=<active> HTTP/1.1
Host: plantdescriptionengine:8432
Authorization: Bearer IDENTITY-TOKEN//<token>
\end{lstlisting}

\begin{lstlisting}[label={lst:query-resp},caption={An \oref{query} response.}]
{
  "entries": [
    {
      "id": 42,
      "plantDescription": "Line 1",
      "active": true,
      "include": [
        1
      ],
      "systems": [
        {
          "systemId": "sensor",
          "systemName": "TemperatureSensor",
          "metadata": {
            "unit": "celsius"
          },
          "ports": [
            {}
          ]
        }
      ],
      "connections": [
        {
          "priority": 1
        }
      ],
      "createdAt": "2026-10-22T12:00:00Z",
      "updatedAt": "2026-10-22T12:00:00Z"
    }
  ],
  "count": 1
}
\end{lstlisting}

\osubsection{get}{GET}{/\allowbreak{}plantdescriptionengine/\allowbreak{}mgmt/\allowbreak{}plant-descriptions/\allowbreak{}\fparam{\{id\}}}{\pref{String}}{\mref{PlantDescriptionResponse}}

Returns a plant description.

Path parameter: \fparam{id}.

A successful call returns \texttt{200 OK} with the output in the body.

\begin{lstlisting}[label={lst:get},caption={An \oref{get} invocation.}]
GET /plantdescriptionengine/mgmt/plant-descriptions/<id> HTTP/1.1
Host: plantdescriptionengine:8432
Authorization: Bearer IDENTITY-TOKEN//<token>
\end{lstlisting}

\begin{lstlisting}[label={lst:get-resp},caption={An \oref{get} response.}]
{
  "id": 42,
  "plantDescription": "Line 1",
  "active": true,
  "include": [
    1
  ],
  "systems": [
    {
      "systemId": "sensor",
      "systemName": "TemperatureSensor",
      "metadata": {
        "unit": "celsius"
      },
      "ports": [
        {
          "portName": "out",
          "serviceDefinition": "temperatureReading",
          "serviceInterface": "string",
          "metadata": {
            "unit": "celsius"
          },
          "consumer": "TemperatureDisplay"
        }
      ]
    }
  ],
  "connections": [
    {
      "consumer": {
        "systemId": "sensor",
        "portName": "out"
      },
      "producer": {
        "systemId": "sensor",
        "portName": "out"
      },
      "priority": 1
    }
  ],
  "createdAt": "2026-10-22T12:00:00Z",
  "updatedAt": "2026-10-22T12:00:00Z"
}
\end{lstlisting}

\osubsection{create}{POST}{/\allowbreak{}plantdescriptionengine/\allowbreak{}mgmt/\allowbreak{}plant-descriptions}{\mref{PlantDescriptionRequest}}{\mref{PlantDescriptionResponse}}

Creates a plant description; if active, it replaces the active one and the rules are synchronized.

A successful call returns \texttt{201 Created} with the output in the body.

\begin{lstlisting}[label={lst:create},caption={An \oref{create} invocation.}]
POST /plantdescriptionengine/mgmt/plant-descriptions HTTP/1.1
Host: plantdescriptionengine:8432
Authorization: Bearer IDENTITY-TOKEN//<token>
Content-Type: application/json

{
  "plantDescription": "Line 1",
  "active": true,
  "include": [
    1
  ],
  "systems": [
    {
      "systemId": "sensor",
      "systemName": "TemperatureSensor",
      "metadata": {
        "unit": "celsius"
      },
      "ports": [
        {
          "portName": "out",
          "serviceDefinition": "temperatureReading",
          "serviceInterface": "string",
          "metadata": {
            "unit": "celsius"
          },
          "consumer": "TemperatureDisplay"
        }
      ]
    }
  ],
  "connections": [
    {
      "consumer": {
        "systemId": "sensor",
        "portName": "out"
      },
      "producer": {
        "systemId": "sensor",
        "portName": "out"
      },
      "priority": 1
    }
  ]
}
\end{lstlisting}

\begin{lstlisting}[label={lst:create-resp},caption={An \oref{create} response.}]
{
  "id": 42,
  "plantDescription": "Line 1",
  "active": true,
  "include": [
    1
  ],
  "systems": [
    {
      "systemId": "sensor",
      "systemName": "TemperatureSensor",
      "metadata": {
        "unit": "celsius"
      },
      "ports": [
        {
          "portName": "out",
          "serviceDefinition": "temperatureReading",
          "serviceInterface": "string",
          "metadata": {
            "unit": "celsius"
          },
          "consumer": "TemperatureDisplay"
        }
      ]
    }
  ],
  "connections": [
    {
      "consumer": {
        "systemId": "sensor",
        "portName": "out"
      },
      "producer": {
        "systemId": "sensor",
        "portName": "out"
      },
      "priority": 1
    }
  ],
  "createdAt": "2026-10-22T12:00:00Z",
  "updatedAt": "2026-10-22T12:00:00Z"
}
\end{lstlisting}

\osubsection{replace}{PUT}{/\allowbreak{}plantdescriptionengine/\allowbreak{}mgmt/\allowbreak{}plant-descriptions/\allowbreak{}\fparam{\{id\}}}{\mref{PlantDescriptionRequest}}{\mref{PlantDescriptionResponse}}

Replaces a plant description.

Path parameter: \fparam{id}.

A successful call returns \texttt{200 OK} with the output in the body.

\begin{lstlisting}[label={lst:replace},caption={An \oref{replace} invocation.}]
PUT /plantdescriptionengine/mgmt/plant-descriptions/<id> HTTP/1.1
Host: plantdescriptionengine:8432
Authorization: Bearer IDENTITY-TOKEN//<token>
Content-Type: application/json

{
  "plantDescription": "Line 1",
  "active": true,
  "include": [
    1
  ],
  "systems": [
    {
      "systemId": "sensor",
      "systemName": "TemperatureSensor",
      "metadata": {
        "unit": "celsius"
      },
      "ports": [
        {
          "portName": "out",
          "serviceDefinition": "temperatureReading",
          "serviceInterface": "string",
          "metadata": {
            "unit": "celsius"
          },
          "consumer": "TemperatureDisplay"
        }
      ]
    }
  ],
  "connections": [
    {
      "consumer": {
        "systemId": "sensor",
        "portName": "out"
      },
      "producer": {
        "systemId": "sensor",
        "portName": "out"
      },
      "priority": 1
    }
  ]
}
\end{lstlisting}

\begin{lstlisting}[label={lst:replace-resp},caption={An \oref{replace} response.}]
{
  "id": 42,
  "plantDescription": "Line 1",
  "active": true,
  "include": [
    1
  ],
  "systems": [
    {
      "systemId": "sensor",
      "systemName": "TemperatureSensor",
      "metadata": {
        "unit": "celsius"
      },
      "ports": [
        {
          "portName": "out",
          "serviceDefinition": "temperatureReading",
          "serviceInterface": "string",
          "metadata": {
            "unit": "celsius"
          },
          "consumer": "TemperatureDisplay"
        }
      ]
    }
  ],
  "connections": [
    {
      "consumer": {
        "systemId": "sensor",
        "portName": "out"
      },
      "producer": {
        "systemId": "sensor",
        "portName": "out"
      },
      "priority": 1
    }
  ],
  "createdAt": "2026-10-22T12:00:00Z",
  "updatedAt": "2026-10-22T12:00:00Z"
}
\end{lstlisting}

\osubsection{update}{PATCH}{/\allowbreak{}plantdescriptionengine/\allowbreak{}mgmt/\allowbreak{}plant-descriptions/\allowbreak{}\fparam{\{id\}}}{\mref{PlantDescriptionUpdate}}{\mref{PlantDescriptionResponse}}

Changes the given fields of a plant description.

Path parameter: \fparam{id}.

A successful call returns \texttt{200 OK} with the output in the body.

\begin{lstlisting}[label={lst:update},caption={An \oref{update} invocation.}]
PATCH /plantdescriptionengine/mgmt/plant-descriptions/<id> HTTP/1.1
Host: plantdescriptionengine:8432
Authorization: Bearer IDENTITY-TOKEN//<token>
Content-Type: application/json

{
  "plantDescription": "Line 1",
  "active": true,
  "include": [
    1
  ],
  "systems": [
    {
      "systemId": "sensor",
      "systemName": "TemperatureSensor",
      "metadata": {
        "unit": "celsius"
      },
      "ports": [
        {
          "portName": "out",
          "serviceDefinition": "temperatureReading",
          "serviceInterface": "string",
          "metadata": {
            "unit": "celsius"
          },
          "consumer": "TemperatureDisplay"
        }
      ]
    }
  ],
  "connections": [
    {
      "consumer": {
        "systemId": "sensor",
        "portName": "out"
      },
      "producer": {
        "systemId": "sensor",
        "portName": "out"
      },
      "priority": 1
    }
  ]
}
\end{lstlisting}

\begin{lstlisting}[label={lst:update-resp},caption={An \oref{update} response.}]
{
  "id": 42,
  "plantDescription": "Line 1",
  "active": true,
  "include": [
    1
  ],
  "systems": [
    {
      "systemId": "sensor",
      "systemName": "TemperatureSensor",
      "metadata": {
        "unit": "celsius"
      },
      "ports": [
        {
          "portName": "out",
          "serviceDefinition": "temperatureReading",
          "serviceInterface": "string",
          "metadata": {
            "unit": "celsius"
          },
          "consumer": "TemperatureDisplay"
        }
      ]
    }
  ],
  "connections": [
    {
      "consumer": {
        "systemId": "sensor",
        "portName": "out"
      },
      "producer": {
        "systemId": "sensor",
        "portName": "out"
      },
      "priority": 1
    }
  ],
  "createdAt": "2026-10-22T12:00:00Z",
  "updatedAt": "2026-10-22T12:00:00Z"
}
\end{lstlisting}

\osubsection{remove}{DELETE}{/\allowbreak{}plantdescriptionengine/\allowbreak{}mgmt/\allowbreak{}plant-descriptions/\allowbreak{}\fparam{\{id\}}}{\pref{String}}{}

Removes a plant description and its rules.

Path parameter: \fparam{id}.

A successful call returns \texttt{204 No Content}.

\begin{lstlisting}[label={lst:remove},caption={An \oref{remove} invocation.}]
DELETE /plantdescriptionengine/mgmt/plant-descriptions/<id> HTTP/1.1
Host: plantdescriptionengine:8432
Authorization: Bearer IDENTITY-TOKEN//<token>
\end{lstlisting}

\newpage

\section{Data Models}
\label{sec:model}

Here, all data objects that can be part of the service calls associated with this service are listed in alphabetic order.
Note that each subsection, which describes one type of object, begins with the \textit{struct} keyword, which is meant to denote a JSON \pref{Object} that must contain certain fields, or names, with values conforming to explicitly named types.
Fields that are not needed may be omitted or \texttt{null} in requests; empty fields are omitted in responses.
As a complement to the primary types defined in this section, there is also a list of secondary types in Section \ref{sec:model:primitives}, which are used to represent things like hashes, identifiers and texts.

\msubsection{struct}{Connection}

A connection from a consumer port to a producer port.


\noindent\begin{tabularx}{\textwidth}{| >{\raggedright\arraybackslash}p{3.9cm} | >{\raggedright\arraybackslash\small}p{5.3cm} | X |} \hline
\rowcolor{gray!33} Object Field & Value Type & Description \\ \hline
"consumer" & \mref{SystemPort} & The consumer port of the connection. \\ \hline
"producer" & \mref{SystemPort} & The producer port of the connection. \\ \hline
"priority" & \pref{Number} & Priority of the resulting orchestration rules (1 highest). \\ \hline
\end{tabularx}



\msubsection{struct}{ErrorMessage}




\noindent\begin{tabularx}{\textwidth}{| >{\raggedright\arraybackslash}p{3.9cm} | >{\raggedright\arraybackslash\small}p{5.3cm} | X |} \hline
\rowcolor{gray!33} Object Field & Value Type & Description \\ \hline
"error\allowbreak{}Message" & \pref{String} & Error text, if the event is an error. \\ \hline
"error\allowbreak{}Code" & \pref{Number} &  \\ \hline
"exception\allowbreak{}Type" & \pref{ExceptionType} &  \\ \hline
"origin" & \pref{String} &  \\ \hline
\end{tabularx}



\msubsection{struct}{PdeSystem}

A system of the plant.


\noindent\begin{tabularx}{\textwidth}{| >{\raggedright\arraybackslash}p{3.9cm} | >{\raggedright\arraybackslash\small}p{5.3cm} | X |} \hline
\rowcolor{gray!33} Object Field & Value Type & Description \\ \hline
"system\allowbreak{}Id" & \pref{String} & Identifier of the system within the plant description. \\ \hline
"system\allowbreak{}Name" & \pref{String} & Name of the registered system (optional; metadata may identify it instead). \\ \hline
"metadata" & \pref{Object}$<$\allowbreak{}\pref{String}$>$ & Metadata the registered system must have. \\ \hline
"ports" & \pref{Array}$<$\allowbreak{}\mref{Port}$>$ & Services the system consumes or produces. \\ \hline
\end{tabularx}



\msubsection{struct}{PlantDescriptionListResponse}

List of plant descriptions.


\noindent\begin{tabularx}{\textwidth}{| >{\raggedright\arraybackslash}p{3.9cm} | >{\raggedright\arraybackslash\small}p{5.3cm} | X |} \hline
\rowcolor{gray!33} Object Field & Value Type & Description \\ \hline
"entries" & \pref{Array}$<$\allowbreak{}\mref{PlantDescriptionResponse}$>$ & Matching records of the current page. \\ \hline
"count" & \pref{Number} & Total number of matching records (all pages). \\ \hline
\end{tabularx}



\msubsection{struct}{PlantDescriptionRequest}

A plant description to create or replace.


\noindent\begin{tabularx}{\textwidth}{| >{\raggedright\arraybackslash}p{3.9cm} | >{\raggedright\arraybackslash\small}p{5.3cm} | X |} \hline
\rowcolor{gray!33} Object Field & Value Type & Description \\ \hline
"plant\allowbreak{}Description" & \pref{String} & Name of the plant description. \\ \hline
"active" & \pref{Boolean} & Make this the active plant description. \\ \hline
"include" & \pref{Array}$<$\allowbreak{}\pref{Number}$>$ & Identifiers of included plant descriptions. \\ \hline
"systems" & \pref{Array}$<$\allowbreak{}\mref{PdeSystem}$>$ & Systems to create or update. \\ \hline
"connections" & \pref{Array}$<$\allowbreak{}\mref{Connection}$>$ & Connections between consumer and producer ports. \\ \hline
\end{tabularx}



\msubsection{struct}{PlantDescriptionResponse}

A stored plant description.


\noindent\begin{tabularx}{\textwidth}{| >{\raggedright\arraybackslash}p{3.9cm} | >{\raggedright\arraybackslash\small}p{5.3cm} | X |} \hline
\rowcolor{gray!33} Object Field & Value Type & Description \\ \hline
"id" & \pref{Number} & Identifier of the plant description. \\ \hline
"plant\allowbreak{}Description" & \pref{String} & Name of the plant description. \\ \hline
"active" & \pref{Boolean} & Whether this is the active plant description. \\ \hline
"include" & \pref{Array}$<$\allowbreak{}\pref{Number}$>$ & Identifiers of included plant descriptions. \\ \hline
"systems" & \pref{Array}$<$\allowbreak{}\mref{PdeSystem}$>$ & Systems to create or update. \\ \hline
"connections" & \pref{Array}$<$\allowbreak{}\mref{Connection}$>$ & Connections between consumer and producer ports. \\ \hline
"created\allowbreak{}At" & \pref{DateTime} & UTC time of creation. \\ \hline
"updated\allowbreak{}At" & \pref{DateTime} & UTC time of the last change. \\ \hline
\end{tabularx}



\msubsection{struct}{PlantDescriptionUpdate}

Fields of a plant description to change; absent fields are kept.


\noindent\begin{tabularx}{\textwidth}{| >{\raggedright\arraybackslash}p{3.9cm} | >{\raggedright\arraybackslash\small}p{5.3cm} | X |} \hline
\rowcolor{gray!33} Object Field & Value Type & Description \\ \hline
"plant\allowbreak{}Description" & \pref{String} & Name of the plant description. \\ \hline
"active" & \pref{Boolean} & Make this the active plant description. \\ \hline
"include" & \pref{Array}$<$\allowbreak{}\pref{Number}$>$ & Identifiers of included plant descriptions. \\ \hline
"systems" & \pref{Array}$<$\allowbreak{}\mref{PdeSystem}$>$ & Systems to create or update. \\ \hline
"connections" & \pref{Array}$<$\allowbreak{}\mref{Connection}$>$ & Connections between consumer and producer ports. \\ \hline
\end{tabularx}



\msubsection{struct}{Port}

A service a plant system consumes or produces.


\noindent\begin{tabularx}{\textwidth}{| >{\raggedright\arraybackslash}p{3.9cm} | >{\raggedright\arraybackslash\small}p{5.3cm} | X |} \hline
\rowcolor{gray!33} Object Field & Value Type & Description \\ \hline
"port\allowbreak{}Name" & \pref{String} & Name of the port within its system. \\ \hline
"service\allowbreak{}Definition" & \pref{String} & Service definition name. \\ \hline
"service\allowbreak{}Interface" & \pref{String} & Required interface template of the port, if any. \\ \hline
"metadata" & \pref{Object}$<$\allowbreak{}\pref{String}$>$ & Metadata the service instance must have. \\ \hline
"consumer" & \pref{Boolean} & True for a consumed service, false for a produced one. \\ \hline
\end{tabularx}



\msubsection{struct}{SystemPort}

Reference to a port of a plant system.


\noindent\begin{tabularx}{\textwidth}{| >{\raggedright\arraybackslash}p{3.9cm} | >{\raggedright\arraybackslash\small}p{5.3cm} | X |} \hline
\rowcolor{gray!33} Object Field & Value Type & Description \\ \hline
"system\allowbreak{}Id" & \pref{String} & Identifier of the system within the plant description. \\ \hline
"port\allowbreak{}Name" & \pref{String} & Name of the port within its system. \\ \hline
\end{tabularx}



\subsection{Primitives}
\label{sec:model:primitives}

As all messages are encoded using the JSON format \cite{rfc8259}, the following primitive constructs, part of that standard, become available.
Note that the official standard is defined in terms of parsing rules, while this list only concerns syntactic information.
Furthermore, the \pref{Object} and \pref{Array} types are given optional generic type parameters, which are used in this document to signify when pair values or elements are expected to conform to certain types.

\noindent\begin{tabularx}{\textwidth}{| p{4.6cm} | X |} \hline
\rowcolor{gray!33} Type & Description \\ \hline
\pdef{Array} & A JSON Array \cite{rfc8259}; \pref{Array}$<$A$>$ has elements of type A. \\ \hline
\pdef{Boolean} & JSON Boolean, \texttt{true} or \texttt{false}. \\ \hline
\pdef{DateTime} & JSON String with an ISO 8601 UTC timestamp, e.g. \texttt{2026-10-22T12:00:00Z}. \\ \hline
\pdef{ExceptionType} & One of \texttt{ARROWHEAD}, \texttt{INVALID\_\allowbreak{}PARAMETER}, \texttt{AUTH}, \texttt{FORBIDDEN}, \texttt{DATA\_\allowbreak{}NOT\_\allowbreak{}FOUND}, \texttt{TIMEOUT}, \texttt{LOCKED}, \texttt{INTERNAL\_\allowbreak{}SERVER\_\allowbreak{}ERROR}, \texttt{EXTERNAL\_\allowbreak{}SERVER\_\allowbreak{}ERROR} (JSON String). \\ \hline
\pdef{Number} & JSON Number, an integer. \\ \hline
\pdef{Object} & JSON Object; \pref{Object}$<$A$>$ has string keys and values of type A. \\ \hline
\pdef{String} & JSON String. \\ \hline
\end{tabularx}

\newpage

\bibliographystyle{IEEEtran}
\bibliography{bibliography}

\newpage

\section{Revision History}
\subsection{Amendments}

\noindent\begin{tabularx}{\textwidth}{| p{1cm} | p{3cm} | p{2cm} | X | p{4cm} |} \hline
\rowcolor{gray!33} No. & Date & Version & Subject of Amendments & Author \\ \hline

1 & 2026-10-23 & \arrowversion & Document created from the Arrowhead LaTeX templates for the 5.2.2 Java Spring implementation, including the local extensions of this workspace. & Jerker Delsing \\ \hline
2 & 2026-10-30 & \arrowversion & Checked against the primary Java source tree (Eclipse-Arrowhead-java): services, operations, paths and data models unchanged. & Claude Opus 5.5 (draft) \\ \hline
\end{tabularx}

\subsection{Quality Assurance}

\noindent\begin{tabularx}{\textwidth}{| p{1cm} | p{3cm} | p{2cm} | X |} \hline
\rowcolor{gray!33} No. & Date & Version & Approved by \\ \hline

1 & & \arrowversion  &  \\ \hline

\end{tabularx}

\end{document}
